%% Section 3: Related Literature

\section{Related Literature}
\label{sec:literature}

This paper contributes to four bodies of work as it relates to the economic effects of
cash supply shocks, the role of digital payment technology in developing
economies, the theory of money and payment systems, and the use of satellite
nighttime lights as a high-frequency economic indicator.

\subsection{Cash Supply Shocks and Currency Policy}

The empirical literature on cash supply shocks is anchored by India's
November 2016 demonetisation, when the government abruptly invalidated
86 percent of the currency in circulation. \citet{chodorow2020cash} exploit
geographic distribution of demonetized versus newly printed 
bank notes across India's districts, and find that districts with fewer
bank branches and higher cash reliance experienced significantly larger
declines in economic activity, measured by nighttime lights and formal
employment. In sub-Saharan Africa specifically in Nigeria, a relevant margin of 
financial access is  digital
payment technology that operates independently of physical branch
infrastructure. This motivates the focus of this study on a fintech-based
treatment variable and isolates a \emph{digital payment substitution channel}.

\citet{lahiri2020demonetization} provides a complementary macroeconomic
accounting of the 2016 episode, documenting that the GDP cost was concentrated
in the informal sector and in the months immediately following the shock.
\citet{rogoff2015costs} argues that the aggregate macroeconomic benefits of
reducing paper currency use are substantial.  Evidence from this current study points to an
important asymmetry where reducing cash dependence through digital payment adoption is
welfare-improving, but an abrupt withdrawal of currency without sufficient
digital infrastructure to substitute imposes large real costs concentrated in
the least financially included segments of the economy.

\subsection{Digital Payments and Financial Inclusion}

The literature has established the welfare effects of mobile money and digital
payment technology in low-income settings. \citet{jack2014risk} provide the
foundational causal evidence from Kenya's M-Pesa system, showing that mobile
money enables households to share risk more effectively and smooth consumption
against income shocks. \citet{suri2016mobile} extend this to long-run poverty
outcomes, finding that mobile money adoption increased per-capita consumption
and reduced poverty, with larger effects for female-headed households.
\citet{MbitiWeil2016} finds that increased use of M-Pesa lowers the propensity to use informal savings 
mechanisms and raises the likelihood of being banked.

In West Africa, \citet{aker2016payment} evaluate a mobile-money
cash transfer programme in Niger and find significant reductions in transaction
costs and improvements in household welfare relative to traditional payment
channels. These studies establish that digital payment access provides direct
welfare benefits through risk-sharing and consumption smoothing. The
contribution of this study is to identify digital payment
infrastructure as a stabiliser that insulates the real economy
from disruptions in the supply of physical currency. This channel operates at
the state level rather than the household level and is distinct from the
individual-level welfare gains that had been shown in earlier studies.

Financial inclusion literature, \citet{beck2007finance},
\citet{allen2016foundations}, \citet{burgess2005rural}, had shown  that access
to formal financial services promotes household welfare and firm productivity
through credit access, risk management, and reduced transaction costs. The
naira crunch episode allows us to estimate the value of financial
infrastructure through its performance under stress, when infrastructure 
serves as a substitute for a disrupted payments medium.

\subsection{Monetary Theory and Payment Technology}

The theoretical positioning for a digital-payment substitution channel can be traced
to the cash-in-advance (\textsc{cia}) tradition originating with
\citet{lucas1987money}. In these models, a subset of transactions must be financed with liquid balances 
held prior to purchase. Consequently, an unexpected reduction in the availability of currency tightens 
liquidity constraints, restricts transactions, and reduces economic activity. 
Building on \citet{Alvarez2009} that show that financial innovation lowers the demand for transaction balances by reducing 
reliance on cash, access to digital payment technologies can be interpreted as relaxing transaction-related 
liquidity constraints. Enabling some transactions to be settled electronically rather than with physical currency, 
digital payment systems reduce dependence on cash as the sole medium of exchange. In the context of a cash shortage, 
agents with greater access to digital payment technologies should 
therefore be better able to maintain transactions and economic activity.

The empirical design of this study operationalises this mechanism at the subnational level.
States with higher pre-shock fintech adoption have, in aggregate, a lower
dependence on physical currency for transaction settlement; when cash supply
contracts, the real \textsc{cia} constraint binds less severely in these
states. Section~\ref{sec:theory} formalises this intuition and derives the
testable predictions that the empirical sections evaluate.

\subsection{Satellite Nighttime Lights as Economic Measurement}

High-frequency satellite data have become a complement to official
statistics where administrative data are unavailable at monthly or subnational
resolution. \citet{henderson2012measuring} establish the foundational result
that nighttime-light radiance reliably proxies economic activity, especially
for countries with weak national accounts. \citet{donaldson2016view} survey
applications in economics ranging from road infrastructure to conflict
detection. \citet{gibson2020night} discuss appropriate uses and limitations,
including sensitivity to cloud cover, overglow, and saturation in dense urban
areas. \citet{bruederle2019nighttime} validate nighttime lights as a proxy for
local human development across Africa, supporting their use in the
sub-Saharan context.

For this paper, a contribution of the \textsc{nasa} \VIIRS{} Black Marble
product is not as a primary outcome but as an identification tool. The 45
pre-shock monthly observations of LGA-level luminosity (January 2019 to
September 2022) permit a test of parallel trends in the state-level
fintech treatment variable, validating the identifying assumption for
the firm-level specifications in Section~\ref{sec:results_micro}.
